\section{Related Work}
\label{sec:related_work}

\paragraph{Automated Media Bias Detection.}
Research on computational media bias detection has grown substantially in recent years \cite{hamborg2019automated, rodrigo2024systematic}.
Early approaches relied on lexical features and sentiment analysis to identify slanted language \cite{recasens2013linguistic}.
More recent work employs transformer-based models for document-level \cite{baly2020detect} and sentence-level classification \cite{lei2022sentence}, often leveraging datasets such as MBIC \cite{spinde2021mbic}, BABE \cite{spinde2021neural}, and BASIL \cite{fan2019plain}.
For Spanish media, the MBBMD dataset \cite{rodrigo2025mbbmd} provides annotated news articles with fine-grained bias categories.
A common limitation across these approaches is their exclusive focus on textual content, ignoring the social context in which news is consumed and discussed.

\paragraph{User Interactions as Credibility Signals.}
Several studies have explored user behavior as a signal for content quality and credibility.
\citet{baly2018predicting} demonstrated that source-level features, including Wikipedia metadata and web traffic patterns, can predict factuality and bias of news outlets.
\citet{popat2018declare} combined article content with user reactions on social media for claim credibility assessment.
In the context of social news aggregation, \citet{glenski2017consumers} analyzed how Reddit users engage with news from sources of varying reliability.
These works establish that user interaction patterns carry information about content properties, but none have specifically investigated the relationship between media bias labels and fine-grained interaction features on a social news platform.

\paragraph{Community Polarization and Echo Chambers.}
The role of platform structure in shaping media consumption has been extensively studied.
\citet{cinelli2021echo} quantified echo chamber effects across multiple social media platforms, finding significant content segregation.
\citet{bakshy2015exposure} showed that both algorithmic curation and individual choice contribute to ideological isolation on Facebook.
On Reddit, \citet{soliman2019characterization} characterized community-level partisanship through cross-posting patterns.
Our work extends this line of research by analyzing whether media outlet communities, detected through shared commenter bases, correlate with systematic differences in bias exposure.

\paragraph{Sentiment and Engagement in News Discussion.}
Prior research has linked news topic and framing to audience emotional responses.
\citet{reis2015breaking} found that different news categories elicit distinct comment sentiment distributions.
\citet{park2016supporting} studied how editorial framing influences reader reactions in online forums.
More recently, \citet{rathje2021out} demonstrated that outrage-inducing content drives higher engagement on social media.
We build on these findings by examining whether the specific dimension of media bias, rather than topic or emotional valence alone, predicts measurable differences in comment sentiment, emotion, and engagement dynamics.
